\paperchapter{December 2002}{1 December 2002}{B}{50}{GSET 2002 (II)}
\keytable{\kcell{1}{C} & \kcell{2}{A} & \kcell{3}{C} & \kcell{4}{C} & \kcell{5}{D} & \kcell{6}{D} & \kcell{7}{B} & \kcell{8}{B} & \kcell{9}{D} & \kcell{10}{B} \\
\kcell{11}{A} & \kcell{12}{C} & \kcell{13}{A} & \kcell{14}{D} & \kcell{15}{C} & \kcell{16}{B} & \kcell{17}{C} & \kcell{18}{A} & \kcell{19}{A} & \kcell{20}{D} \\
\kcell{21}{B} & \kcell{22}{C} & \kcell{23}{C} & \kcell{24}{B} & \kcell{25}{D} & \kcell{26}{C} & \kcell{27}{C} & \kcell{28}{A} & \kcell{29}{D} & \kcell{30}{B} \\
\kcell{31}{A} & \kcell{32}{B} & \kcell{33}{D} & \kcell{34}{C} & \kcell{35}{B} & \kcell{36}{C} & \kcell{37}{D} & \kcell{38}{A} & \kcell{39}{A} & \kcell{40}{C} \\
\kcell{41}{B} & \kcell{42}{D} & \kcell{43}{A} & \kcell{44}{B} & \kcell{45}{D} & \kcell{46}{B} & \kcell{47}{D} & \kcell{48}{C} & \kcell{49}{A} & \kcell{50}{B}}
\begin{question}{1}
{Complete the series: 1, 3, 4, 8, 16, ?\par}
\optsfour{30}{31}{32}{33}
\answer{C}{32}
\begin{explain}
Each term (from the third) is the \textbf{sum of all previous terms}: $1+3=4$, $1+3+4=8$, $1+3+4+8=16$, so next $=1+3+4+8+16=\mathbf{32}$.
\end{explain}
\end{question}

\begin{question}{2}
{Complete the series: 10, 12, 16, 22, 30, ?\par}
\optsfour{40}{50}{60}{70}
\answer{A}{40}
\begin{explain}
The differences are $2, 4, 6, 8$, so next $+10$: $30+10=\mathbf{40}$.
\end{explain}
\end{question}

\begin{question}{3}
{Complete the series: 240, 240, 120, 40, ?\par}
\optsfour{30}{20}{10}{0}
\answer{C}{10}
\begin{explain}
Divide successively by 1, 2, 3, 4: $240, 240, 120, 40, \mathbf{10}$.
\end{explain}
\end{question}

\begin{question}{4}
{If CENTRAL has been coded as XVMGIZO, what will be the correct code for PHYSICS?\par}
\optsfour{KSBRHHX}{HXBHKRX}{KSBHRXH}{HXRHSBK}
\answer{C}{KSBHRXH}
\begin{explain}
CENTRAL $\rightarrow$ XVMGIZO uses the \textbf{reverse-alphabet} code (A$\leftrightarrow$Z, B$\leftrightarrow$Y, C$\leftrightarrow$X, \ldots; the two letters add to 27).
PHYSICS: P$\rightarrow$K, H$\rightarrow$S, Y$\rightarrow$B, S$\rightarrow$H, I$\rightarrow$R, C$\rightarrow$X, S$\rightarrow$H, giving \textbf{KSBHRXH}.
\end{explain}
\end{question}

\begin{question}{5}
{If GOAT is coded as HPBU, what will be the correct code for LIME ?\par}
\optsfour{ORNZ}{MJND}{KHLF}{MJNF}
\answer{D}{MJNF}
\begin{explain}
Each letter moves one step forward (G$\rightarrow$H, O$\rightarrow$P, \ldots). LIME $\rightarrow$ \textbf{MJNF}.
\end{explain}
\end{question}

\begin{question}{6}
{If CHEMISTRY has been coded as DGFLJRUQZ, what will be the correct code for DESIGN?\par}
\optsfour{EDTHMN}{EDTMHN}{EDTHNN}{EDTHHM}
\answer{D}{EDTHHM}
\begin{explain}
CHEMISTRY $\rightarrow$ DGFLJRUQZ: the shifts alternate $+1, -1$.
DESIGN: D$+1=$E, E$-1=$D, S$+1=$T, I$-1=$H, G$+1=$H, N$-1=$M, giving \textbf{EDTHHM}.
\end{explain}
\end{question}

\begin{question}{7}
{If CARPET is coded as XZIKVG, what will be the correct code for MARKET?\par}
\optsfour{NZIPUG}{NZIPVG}{NZIQVG}{NZIPTG}
\answer{B}{NZIPVG}
\begin{explain}
CARPET $\rightarrow$ XZIKVG is again the reverse-alphabet code. MARKET $\rightarrow$ N, Z, I, P, V, G $=$ \textbf{NZIPVG}.
\end{explain}
\end{question}

\begin{question}{8}
{The quality of research is judged by the:\par}
\optsone{Relevance of research}{Methodology followed in conducting the research}{Depth of the research}{Experience of the researcher}
\answer{B}{Methodology followed in conducting the research}
\begin{explain}
The quality of research is judged mainly by the \textbf{methodology} followed: the rigour of design, sampling, tools and analysis determines the validity of the findings.
\end{explain}
\end{question}

\begin{question}{9}
{Which one of the following statements is not correct?\par}
\optstwo{Any person can do research.}{Research may be without a hypothesis.}{Research requires library facilities.}{Research is not always systematic.}
\answer{D}{Research is not always systematic.}
\begin{explain}
Research is by definition a \emph{systematic} inquiry. So ``research is not always systematic'' is \textbf{not correct}.
\end{explain}
\end{question}

\begin{question}{10}
{The first step of research is:\par}
\optstwo{Selection of Problem}{Identification of Problem}{Definition of Problem}{Locating the Problem}
\answer{B}{Identification of Problem}
\begin{explain}
The first step is \textbf{identification of the problem}; selection and definition follow.
\end{explain}
\end{question}

\begin{question}{11}
{A good researcher must possess:\par}
\optstwo{an inquiry mind}{a master's degree}{a good health}{enough funds}
\answer{A}{an inquiry mind}
\begin{explain}
A good researcher must above all possess an \textbf{inquiring mind} (curiosity).
\end{explain}
\end{question}

\begin{question}{12}
{Which one of the following statements is not correct?\par}
\optsone{Hypothesis and assumptions are not same.}{Problem refers to variable.}{Quantitative data can be converted into qualitative data.}{Qualitative data may be converted into quantitative data.}
\answer{C}{Quantitative data can be converted into qualitative data.}
\begin{explain}
Per the key, ``quantitative data can be converted into qualitative data'' is treated as \textbf{not correct}. Statement (D), converting qualitative data into quantitative data by coding and counting, is common practice.
\end{explain}
\end{question}

\begin{context}
\textbf{Read the following passage and answer the questions 13--16.}\par\smallskip
Both plants and animals of many sorts show remarkable changes in form, structure, growth habits, and even mode of reproduction in becoming adapted to different climatic environment, types of food supply, or mode of living. This divergence in response to evolution is commonly expressed by altering the form and function of some part or parts of the organism, the original identity of which is clearly discernible. For example, the creeping foot of the snail is seen in related marine pteropods to be modified into a flapping organ useful for swimming, and is changed into prehensile arms that bear suctorial disks in the squids and other cephalopods. The limbs of various mammals are modified according to several different modes of life-for swift running (cursorial) as in the horse and antelope, for swinging in trees (arboreal) as in the monkeys, for digging (fossorial) as in the moles and gophers, for flying (volant) as in the bats, for swimming (aquatic) as in the seals, whales and dolphins, and for other adaptations. The structures or organs that show main change in connection with this adaptive divergence are commonly identified readily as homologous, in spite of great alterations. Thus, the finger and wristbones of a bat and whale, for instance, have virtually nothing in common except that they are definitely equivalent elements of the mammalian limb.
\end{context}
\begin{question}{13}
{Which of the following is the most appropriate title for the passage, based on its content ?\par}
\optstwo{Adaptive Divergence}{Evolution}{Changes in Organs}{Our Changing Bodies}
\answer{A}{Adaptive Divergence}
\begin{explain}
The passage describes how organisms diverge in form in \emph{adapting} to different environments: \textbf{Adaptive Divergence} (the term the passage itself uses).
\end{explain}
\end{question}

\begin{question}{14}
{The author provides information that would answer which of the following questions?\newline (I) What factors cause change in organisms ?\newline (II) What is the theory of evolution?\newline (III) How are horses' legs related to seals' flippers ?\par}
\optsfour{I only}{II only}{I and II only}{I and III only}
\answer{D}{I and III only}
\begin{explain}
The passage explains that climate, food supply and mode of living cause change (I), and it relates the limbs of horses (running) and seals (swimming) as homologous (III). It does not explain the theory of evolution (II). \textbf{I and III only}.
\end{explain}
\end{question}

\begin{question}{15}
{Which of the following words could best be substituted for 'homologous' without substantially changing the author's meaning?\par}
\optsfour{Altered}{Mammalian}{Corresponding}{Divergent}
\answer{C}{Corresponding}
\begin{explain}
\textbf{Homologous} organs are equivalent or \textbf{corresponding} parts that share an origin despite different functions.
\end{explain}
\end{question}

\begin{question}{16}
{The author's style can best be described as:\par}
\optsfour{Humorous}{Objective}{Esoteric}{Patronizing}
\answer{B}{Objective}
\begin{explain}
The author presents facts and examples without personal bias: an \textbf{objective} style.
\end{explain}
\end{question}

\begin{context}
\textbf{Read the following passage and answer the questions 17--22.}\par\smallskip
The atmosphere is a mixture of several gases. There are about ten chemical elements which remain permanently in gaseous form in the atmosphere under all natural conditions. Of these permanent gases, oxygen makes up about 21 percent and nitrogen about 78 percent. Several other gases, such as argon, carbon dioxide, hydrogen, neon, krypton and Xenon, comprise the remaining 1 percent of the volume of dry air. The amount of water vapour, and its variations in amount and distribution, are of extraordinary importance in weather changes. Atmospheric gases hold in suspension great quantities of dust, pollen, smoke and other impurities which are always present in considerable, but variable amounts.

The atmosphere has no definite upper limits but gradually thins until it becomes imperceptible. Until recently it was assumed that the air above the first few miles gradually grew thinner and colder at a constant rate. It was also assumed that upper air had little influence on weather changes. Recent studies of the upper atmosphere, currently being conducted by earth satellites and missile probings, have shown these assumptions to be incorrect. The atmosphere has three well-defined strata.

The layer of the air next to the earth, which extends upward for about 10 miles, is known as troposphere. On the whole, it makes up about 75 percent of all the weight of the atmosphere. It is the warmest part of the atmosphere because most of the solar radiation is absorbed by the earth's surface, which warms the air immediately surrounding it. A steady decrease of temperature with increasing elevation is a most striking characteristic. The upper layers are colder because of their great distance from the earth's surface and rapid radiation of heat into space. The temperatures within the troposphere decrease about 3.5 degrees per 1,000 foot increase in altitude. Within the troposphere, winds and air currents distribute heat and moisture. Strong winds, called jet streams, are located at the upper levels of the troposphere. These jet streams are both complex and wide spread in occurrence. They normally show a wave-shaped pattern and move from west to east at velocities of 150 mph, but velocities as high as 400 mph have been noted. The influences of changing locations and strengths of jet streams upon weather conditions and patterns are no doubt considerable. Current intensive research may eventually reveal their true significance.

Above the troposphere to a height of about 50 miles is a zone called the stratosphere. The stratosphere is separated from the troposphere by a zone of uniform temperatures called tropopause. Within the lower portions of the stratosphere is a layer of ozone gases which filters outmost of the ultraviolet rays from the sun. The ozone layer varies with air pressure. If this zone was not there, the full blast of the sun's ultraviolet light would burn our skin, blind our eyes, and eventually result in our destruction. Within the stratosphere, the temperature and atmospheric composition are relatively uniform.

The layer upward of about 50 miles is the most fascinating but the least known of these three strata. It is called ionosphere because it consists of electrically charged particles called ions, thrown from the sun. The northern lights (aurora borealis) originate within this highly charged portion of the atmosphere. Its effect upon weather conditions, if any, is as yet unknown.
\end{context}
\begin{question}{17}
{Which of the following titles best summarizes the content of the passage?\par}
\optsone{New methods for calculating the composition of the Atmosphere}{New Evidence concerning the stratification of the Atmosphere}{The Atmosphere: Its nature and importance to our weather}{The underlying causes of Atmospheric Turbulence}
\answer{C}{The Atmosphere: Its nature and importance to our weather}
\begin{explain}
The passage describes the composition and layers of the atmosphere and repeatedly relates them to weather: \textbf{The Atmosphere: its nature and importance to our weather}.
\end{explain}
\end{question}

\begin{question}{18}
{According to the passage, life as we know it exists on the earth because the atmosphere:\par}
\optsone{contains a layer of ozone gases.}{contains electrically charged particles.}{is warmest at the bottom.}{carries the ultraviolet rays of the sun.}
\answer{A}{contains a layer of ozone gases.}
\begin{explain}
The \textbf{ozone layer} filters the Sun's ultraviolet rays; without it ``the full blast of the sun's ultraviolet light would \ldots\ result in our destruction''.
\end{explain}
\end{question}

\begin{question}{19}
{The passage supplies information that would answer which of the following questions ?\newline (I) How do the troposphere and the stratosphere differ ?\newline (II) How does the ionosphere affect the weather?\newline (III) How do earth satellites study the atmosphere ?\par}
\optsfour{I only}{II only}{III only}{I, II and III}
\answer{A}{I only}
\begin{explain}
Only question I is answered (the troposphere and stratosphere are contrasted in temperature and composition). The ionosphere's effect on weather is ``as yet unknown'' (II), and how satellites study the atmosphere is not explained (III). \textbf{I only}.
\end{explain}
\end{question}

\begin{question}{20}
{The passage states that the troposphere is the warmest part of the atmosphere because it :\par}
\optsone{is the closest to the sun.}{radiates heat into space.}{has winds and air current that distribute the heat.}{is warmed by the earth's heat.}
\answer{D}{is warmed by the earth's heat.}
\begin{explain}
The troposphere is warmest ``because most of the solar radiation is absorbed by the earth's surface, which warms the air immediately surrounding it'', i.e.\ it is \textbf{warmed by the earth's heat}.
\end{explain}
\end{question}

\begin{question}{21}
{According to the passage, the atmosphere consists of all of the following, except:\par}
\optstwo{21 percent oxygen}{a definite amount of water vapour}{ten permanent elements}{considerable waste products}
\answer{B}{a definite amount of water vapour}
\begin{explain}
The passage stresses the \emph{variations} in the amount of water vapour, so ``a \textbf{definite} amount of water vapour'' is the exception.
\end{explain}
\end{question}

\begin{question}{22}
{It can be inferred from the passage that at the top of Jungfrau, which is 12,000 feet above the town of Interlaken in Switzerland, the temperature is usually :\par}
\optsone{below freezing}{warmer than in Interlaken}{about 42 degrees colder than on the ground}{about 75 degrees colder than in Interlaken}
\answer{C}{about 42 degrees colder than on the ground}
\begin{explain}
Temperature falls about $3.5^\circ$ per 1,000 ft, so at 12,000 ft: $12\times3.5=\mathbf{42}^\circ$ colder than on the ground.
\end{explain}
\end{question}

\begin{question}{23}
{If I say 'sky is down and winter is hot', I would call a criminal:\par}
\optsfour{Gangster}{A cloud}{Saint}{Delinquent}
\answer{C}{Saint}
\begin{explain}
The speaker uses words in their \emph{opposite} sense (sky is ``down'', winter is ``hot''), so a criminal would be called a \textbf{saint}.
\end{explain}
\end{question}

\begin{question}{24}
{If Dolly's mother is John's father's sister, what relation is John to Dolly's father ?\par}
\optsfour{Cousin}{Nephew}{Uncle}{Brother}
\answer{B}{Nephew}
\begin{explain}
Dolly's mother is John's paternal aunt, so Dolly's father is the husband of John's aunt, i.e.\ John's uncle (by marriage). John is his \textbf{nephew}.
\end{explain}
\end{question}

\begin{question}{25}
{Establish the relationship between the first two items. Find out the same relationship among the choices which follow the pattern of the two items: 20:21:: ?\par}
\optsfour{6:10}{10:B}{S: V}{A: B}
\answer{D}{A: B}
\begin{explain}
20 : 21 are consecutive numbers, corresponding to T : U, consecutive letters. Similarly \textbf{A : B} are consecutive (positions 1, 2).
\end{explain}
\end{question}

\begin{context}
\textbf{Study the pie-chart and answer the questions 26 to 30.} (The percentages below are those implied by the questions and official answers; the original chart was printed without values.)
\begin{center}
\begin{tikzpicture}\pie[radius=1.4,text=pin,font=\scriptsize,color={brand!35,brand!10,brand!25,brand!45,brand!5}]{25/Engineering,12.5/Medical,37.5/Arts,12.5/Commerce,12.5/Science}\end{tikzpicture}
\end{center}
\end{context}
\begin{question}{26}
{What is the percentage of total students enrolled in Medical and Engineering faculties ?\par}
\optsfour{33.50}{50.50}{37.50}{25.50}
\answer{C}{37.50}
\begin{explain}
Medical $+$ Engineering $=12.5\%+25\%=\mathbf{37.5\%}$.
\end{explain}
\end{question}

\begin{question}{27}
{What is the percentage of students enrolled in Arts and Science faculties?\par}
\optsfour{90}{60}{50}{45}
\answer{C}{50}
\begin{explain}
Arts $+$ Science $=37.5\%+12.5\%=\mathbf{50\%}$.
\end{explain}
\end{question}

\begin{question}{28}
{What is the percentage of students enrolled in Arts and Engineering faculties ?\par}
\optsfour{62.50}{75.00}{50.50}{33.33}
\answer{A}{62.50}
\begin{explain}
Arts $+$ Engineering $=37.5\%+25\%=\mathbf{62.5\%}$.
\end{explain}
\end{question}

\begin{question}{29}
{Suppose that the number of students enrolled in Engineering is 320; what is the number of students enrolled in Arts faculty?\par}
\optsfour{450}{460}{470}{480}
\answer{D}{480}
\begin{explain}
Engineering $=25\%=320$, so $1\%=12.8$ and Arts $=37.5\%=37.5\times12.8=\mathbf{480}$.
\end{explain}
\end{question}

\begin{question}{30}
{Suppose, the number of students enrolled in Engineering is 320; what is the total number of students enrolled in Commerce, Arts and Medical faculties?\par}
\optsfour{700}{800}{900}{1000}
\answer{B}{800}
\begin{explain}
Commerce $+$ Arts $+$ Medical $=12.5+37.5+12.5=62.5\%$ of $1280=\mathbf{800}$.
\end{explain}
\end{question}

\begin{question}{31}
{All streets are watches. All watches are eagles. Therefore :\newline (I) All streets are eagles.\newline (II) All watches are streets.\par}
\optstwo{if only inference (I) follows}{if only inference (II) follows}{if both of them follow}{if neither of them follows}
\answer{A}{if only inference (I) follows}
\begin{explain}
\begin{center}\begin{tikzpicture}[scale=0.55,font=\scriptsize]\draw (0,0) circle(1.8); \node at (0,1.35){Eagles}; \draw (0,-0.3) circle(1.1); \node at (0,0.4){Watches}; \draw[fill=brand!15] (0,-0.6) circle(0.55) node{Streets};\end{tikzpicture}\end{center}
All streets are watches and all watches are eagles, so \textbf{all streets are eagles} (I follows). II (all watches are streets) is the illicit converse. \textbf{Only I}.
\end{explain}
\end{question}

\begin{question}{32}
{All boys are mothers. All mothers are father. Therefore :\newline (I) All mothers are boys.\newline (II) All boys are fathers.\par}
\optstwo{if only inference (I) follows}{if only inference (II) follows}{if both of them follow}{if neither of them follows}
\answer{B}{if only inference (II) follows}
\begin{explain}
All boys are mothers and all mothers are fathers, so \textbf{all boys are fathers} (II follows). I (all mothers are boys) is an illicit converse. \textbf{Only II}.
\end{explain}
\end{question}

\begin{question}{33}
{Monkey is a bird. Some birds are clouds. Therefore :\newline (I) Monkey is a cloud.\newline (II) Some clouds are birds.\par}
\optstwo{if only inference (I) follows}{if only inference (II) follows}{if both of them follow}{if neither of them follows}
\answer{D}{if neither of them follows}
\keynote{Strictly, ``Some birds are clouds'' converts validly to ``Some clouds are birds'', so II follows. The official key gives (D).}
\begin{explain}
``Monkey is a bird'' and ``some birds are clouds'' do not show that the monkey is among those birds that are clouds, so I does not follow. II is the simple converse of the second premise. The official key, however, gives \textbf{neither follows}.
\end{explain}
\end{question}

\begin{context}
Expected (E) as well as actual (A) increase in export of different items are given in the following table (values in crores). Take into account the proportions of expected and actual export.
\begin{center}\small
\begin{tabular}{l*{8}{c}}\toprule
& \multicolumn{2}{c}{1990} & \multicolumn{2}{c}{1991} & \multicolumn{2}{c}{1992} & \multicolumn{2}{c}{1993}\\\cmidrule(lr){2-3}\cmidrule(lr){4-5}\cmidrule(lr){6-7}\cmidrule(lr){8-9}
\textbf{Items} & E & A & E & A & E & A & E & A\\\midrule
Rice & 90 & 60 & 120 & 60 & 150 & 60 & 200 & 60\\ Sugar & 100 & 70 & 200 & 140 & 300 & 210 & 400 & 280\\ Wheat & 100 & 75 & 200 & 150 & 300 & 200 & 400 & 300\\ Textile & 500 & 450 & 700 & 600 & 800 & 700 & 1000 & 850\\ Machinery & 1000 & 600 & 1500 & 700 & 2000 & 1500 & 2500 & 1500\\ Electronic goods & 1500 & 1200 & 2000 & 1500 & 2500 & 2000 & 3000 & 2000\\\bottomrule
\end{tabular}\\[4pt]
\textbf{Actual as \% of expected:}\\[2pt]
\begin{tabular}{lcccc}\toprule & 1990 & 1991 & 1992 & 1993\\\midrule
Rice & 66.7 & 50.0 & 40.0 & 30.0\\ Sugar & 70.0 & 70.0 & 70.0 & 70.0\\ Wheat & 75.0 & 75.0 & 66.7 & 75.0\\ Textile & 90.0 & 85.7 & 87.5 & 85.0\\ Machinery & 60.0 & 46.7 & 75.0 & 60.0\\ Electronic goods & 80.0 & 75.0 & 80.0 & 66.7\\\bottomrule
\end{tabular}
\end{center}
\end{context}
\begin{question}{34}
{Which item indicates consistent decrease in the export during the four years ?\par}
\optstwo{Electronic goods}{Textile}{Rice}{Sugar}
\answer{C}{Rice}
\begin{explain}
Only \textbf{rice} falls every year: 66.7\% $\rightarrow$ 50\% $\rightarrow$ 40\% $\rightarrow$ 30\%.
\end{explain}
\end{question}

\begin{question}{35}
{During 1990 which one of the four items had the highest percentage of export (Actual/Expected) ?\par}
\optsfour{Machinery}{Textile}{Wheat}{Sugar}
\answer{B}{Textile}
\begin{explain}
In 1990: Machinery 60\%, \textbf{Textile 90\%}, Wheat 75\%, Sugar 70\%. Textile is highest.
\end{explain}
\end{question}

\begin{question}{36}
{During which year the percent export of wheat was the lowest ?\par}
\optsfour{1990}{1991}{1992}{1993}
\answer{C}{1992}
\begin{explain}
Wheat: 75, 75, \textbf{66.7}, 75. The lowest is \textbf{1992}.
\end{explain}
\end{question}

\begin{question}{37}
{In which year sugar export was the lowest ?\par}
\optsfour{1990}{1992}{1993}{None of these}
\answer{D}{None of these}
\begin{explain}
Sugar's percentage is 70\% in every year, so there is no single lowest year: \textbf{none of these}.
\end{explain}
\end{question}

\begin{question}{38}
{During 1993 which item has the lowest percentage of export ?\par}
\optsfour{Rice}{Sugar}{Wheat}{Machinery}
\answer{A}{Rice}
\begin{explain}
1993: \textbf{Rice 30\%}, Sugar 70\%, Wheat 75\%, Machinery 60\%. Rice is lowest.
\end{explain}
\end{question}

\begin{question}{39}
{Which word does not belong to the other three ?\par}
\optsfour{Think}{See}{Hear}{Smell}
\answer{A}{Think}
\begin{explain}
See, hear and smell are \emph{sense perceptions}; \textbf{think} is a mental (cognitive) process.
\end{explain}
\end{question}

\begin{question}{40}
{Find one group of letters which is different from the other three:\par}
\optsfour{TRP}{SQO}{TVX}{YWU}
\answer{C}{TVX}
\begin{explain}
TRP, SQO and YWU move backwards in steps of 2; \textbf{TVX} moves forwards, so it is different.
\end{explain}
\end{question}

\begin{question}{41}
{Teaching refers to:\par}
\optsone{giving information related to the subject.}{creating environment that is conducive for learning.}{dictating notes in the classroom.}{solving difficulties of students related to the subject.}
\answer{B}{creating environment that is conducive for learning.}
\begin{explain}
Teaching means \textbf{creating an environment conducive to learning}, not merely giving information or dictating notes.
\end{explain}
\end{question}

\begin{question}{42}
{Which one of the following statements in not true?\par}
\optsone{Majority of teachers use lecture method.}{The lecture method is less structured.}{Lecture method is the best for giving more information in less time.}{There are well-defined steps of lecture method.}
\answer{D}{There are well-defined steps of lecture method.}
\begin{explain}
Per the key, the untrue statement is (D). The lecture method is widely used, fairly structured and efficient for conveying information; statement (D) overstates its standing.
\end{explain}
\end{question}

\begin{question}{43}
{Good teaching is best reflected by:\par}
\optsone{meaningful questions asked by students.}{attendance of students.}{number of distinctions.}{pin-drop silence in the class.}
\answer{A}{meaningful questions asked by students.}
\begin{explain}
Good teaching is reflected in the \textbf{meaningful questions asked by students}, which show understanding and curiosity.
\end{explain}
\end{question}

\begin{question}{44}
{The main objective of teaching at higher education level is:\par}
\optsone{to give information.}{to develop reasoning and thinking abilities.}{to develop expression power.}{to develop study habits.}
\answer{B}{to develop reasoning and thinking abilities.}
\begin{explain}
The main objective of higher education teaching is to \textbf{develop reasoning and thinking abilities}.
\end{explain}
\end{question}

\begin{question}{45}
{Noise in the classroom is best indicative of:\par}
\optsone{Only interaction among students.}{Only interaction among teacher and students.}{Only interaction between teacher and student.}{(A) and/or (B).}
\answer{D}{(A) and/or (B).}
\begin{explain}
Noise in class may come from interaction among students \emph{and/or} between teacher and students: \textbf{(A) and/or (B)}.
\end{explain}
\end{question}

\begin{question}{46}
{A student is reading a book at home. In this situation :\par}
\optsone{There is no communication.}{There is a communication.}{There is a two-way communication.}{There is only non-verbal communication.}
\answer{B}{There is a communication.}
\begin{explain}
Reading a book is communication between the author (sender) and the reader (receiver) through the printed text (channel): \textbf{there is communication} (one-way).
\end{explain}
\end{question}

\begin{question}{47}
{The facial expressions of students relate to which element of the communication process?\par}
\optsfour{Message}{Receiver}{Sender}{Channel}
\answer{D}{Channel}
\begin{explain}
Facial expressions are a \textbf{non-verbal channel} through which students send messages (interest, confusion, boredom) back to the teacher. The key therefore relates them to the \textbf{channel} of communication.
\end{explain}
\end{question}

\begin{question}{48}
{Communication refers to:\par}
\optsone{exchange of information only.}{exchange of feelings only.}{exchange of information, feelings, emotions and views.}{exchange of emotions only.}
\answer{C}{exchange of information, feelings, emotions and views.}
\begin{explain}
Communication is the \textbf{exchange of information, feelings, emotions and views}.
\end{explain}
\end{question}

\begin{question}{49}
{The Indian classrooms are generally best characterised by:\par}
\optstwo{One-way verbal communication}{Two-way verbal communication}{One-way communication}{Verbal communication}
\answer{A}{One-way verbal communication}
\begin{explain}
Indian classrooms are typically dominated by the lecture: \textbf{one-way verbal communication}.
\end{explain}
\end{question}

\begin{question}{50}
{Two persons sitting facing each other are engrossed in thinking. In this situation:\par}
\optsone{There is non-verbal communication.}{There is no communication.}{There is two-way non-verbal communication.}{There is less verbal communication.}
\answer{B}{There is no communication.}
\begin{explain}
Two people engrossed in their own thoughts are not exchanging any message, so there is \textbf{no communication} between them.
\end{explain}
\end{question}
